\section{Introduction}

Knowledge distillation~\cite{hinton2015distilling} is a promising approach for transferring the capabilities of large language models (LLMs) to smaller, more efficient models with lower inference cost and improved deployability. Traditional distillation relies on off-policy teacher data, training students with supervised loss or forward KL divergence~\citep{kim2016sequence}. However, this introduces a distribution mismatch between training sequences and those generated by the student at inference time.

On-policy distillation addresses this by having the student generate samples that are corrected by the teacher, typically via reverse KL divergence~\citep{gu2024minillm, agarwal2024policy, lu2025onpolicydistillation}.
As noted by \citet{lu2025onpolicydistillation}, this objective can be seamlessly integrated into standard reinforcement learning~(RL) pipelines. Furthermore, the approach is highly efficient: recent work~\citep{lu2025onpolicydistillation, yang2025qwen3} demonstrates that on-policy distillation can match RL-trained models on math reasoning benchmarks at 10$\times$ lower compute cost than GRPO~\citep{shao2024deepseekmath}, which has made on-policy distillation a standard recipe for distillation. 
% Incorporating RL into this framework further improves efficiency, requiring fewer rollouts to achieve strong performance. 
% Recent work has shown that on-policy distillation can match RL-trained models on math reasoning benchmarks at 10–30× lower compute cost~\citep{lu2025onpolicydistillation}.
%(see \autoref{fig:intro_example}).

% However, reverse KL is a \textit{mode-seeking} objective: while it effectively captures the teacher's dominant modes, we show that it reduces student diversity and yields unstable learning signals at positions where the teacher distribution has high entropy (see supporting results in \S\ref{sec:limit_opd}). This limits the student's ability to preserve the teacher's distributional structure when probability mass is spread across multiple tokens. It is particularly problematic in reasoning tasks, where high-entropy tokens often represent key decision points with multiple valid paths~\cite{wang2025beyond}.


% In this work, we propose \methodfullname~(\method). Our key insight is that reverse KL and forward KL are complementary: reverse KL enables efficient learning on confident predictions, while forward KL's mode-covering property transfers uncertainty and global structure. However, naively applying forward KL degrades efficiency by forcing the student to cover low-probability tails~\cite{gu2023minillm, cha2025knowledge}. We therefore propose an entropy-aware strategy that selectively activates forward KL based on the teacher's token-level uncertainty. Specifically, we use reverse KL in low-entropy regions for efficiency, and forward KL in high-entropy regions to preserve diversity. As a result, the student maintains probability mass across multiple plausible continuations and more faithfully reflects the teacher's underlying structure.

% Empirically, this guidance maintains substantially higher generation diversity, preserving teacher-like uncertainty by retaining significantly more probability mass in high-entropy regions than standard on-policy distillation. This improvement translates into consistent downstream gains: averaged over six mathematical reasoning benchmarks, EOPD improves Avg@8 accuracy by +1.16 and Pass@8 by +1.37 for Qwen3-0.6B-Base, by +0.99 Avg@8 and +2.39 Pass@8 for Qwen3-1.7B-Base, and by +1.80 Avg@8 and +5.05 Pass@8 for Qwen3-4B-Base relative to standard on-policy distillation.

% \yj{add 1~2 sentences in intro }


However, reverse KL is a \textit{mode-seeking} objective: while it effectively captures the teacher's dominant modes, we find that it reduces student diversity and yields unstable learning signals at positions where the teacher distribution has high entropy. This limits the student's ability to preserve the teacher's distributional structure when probability mass is spread across multiple tokens. It is particularly problematic in reasoning tasks, where high-entropy tokens often represent key decision points with multiple valid paths~\cite{wang2025beyond, cheng2025reasoning}. The student then converges to a narrower behavioral repertoire than the teacher, even when reverse KL itself is well optimized.
\vspace{-8pt}
\paragraph{Contribution.}
To address this limitation, we propose \methodfullname~(\method), a distillation framework that balances efficiency and diversity. Our key insight is that reverse KL and forward KL are complementary: reverse KL enables efficient learning on confident predictions, while forward KL's mode-covering property transfers uncertainty and global structure. Our specific contributions are as follows:

\begin{itemize}[leftmargin=1.2em, itemsep=1mm, topsep=1pt]
    \item \textbf{Analysis of Diversity Degradation and Training Instability.}
    We conduct a systematic analysis of token-level entropy (\S\ref{sec:diversity}), revealing that standard on-policy distillation causes diversity collapse, retaining only 6.8\% of high-entropy tokens compared to 18.5\% in the teacher. Furthermore, through a controlled toy experiment (\S\ref{sec:instability}), we demonstrate that the reverse KL objective provides unstable gradient signals when the teacher is uncertain, preventing proper convergence.

    \item \textbf{Entropy-Aware On-Policy Distillation (\method).} We introduce an entropy-aware strategy that dynamically adapts the training objective~(\S\ref{sec:main_method}). By selectively applying reverse KL in low-entropy regions for efficiency and forward KL in high-entropy regions to preserve diversity, \method~effectively transfers the teacher's uncertainty without the computational overhead of naive forward KL.

    \item \textbf{Improvements on Reasoning Benchmarks.} Empirically, \method~maintains substantially higher generation diversity, preserving the teacher's uncertainty by retaining more probability mass in high-entropy regions than standard on-policy distillation~(\S\ref{sec:experiments}). This improvement translates into consistent downstream gains: averaged over six mathematical reasoning benchmarks, \method~improves Avg@8 accuracy by +1.16 and Pass@8 by +1.37 for Qwen3-0.6B-Base, with larger gains of +0.99/+2.39 for the 1.7B model and +1.80/+5.05 for the 4B model.
\end{itemize}

%%%%% prev version
% Knowledge distillation is a widely used paradigm for transferring the capabilities of large, highly expressive models to smaller and more efficient ones, enabling practical deployment without sacrificing performance. While classical distillation methods rely on supervised matching of output distributions, recent work has proposed on-policy distillation, a reinforcement learning–based approach that more directly aligns the student’s generation behavior using the guidance of a stronger teacher. In on-policy distillation, the student model is trained using a dense reward derived from the reverse Kullback–Leibler (KL) divergence between the teacher and student token-level distributions along trajectories sampled from the student policy. This formulation provides informative, fine-grained feedback during generation and has been shown to substantially accelerate learning, allowing student models to approach the performance of reinforcement-learning–trained reasoning models in only a few optimization steps (see \autoref{fig:intro_example}).


% Paragraph 1: LLMs + RL (RLVR) & Need for Small Models
% Paragraph 2: Distillation Research & Emphasizing On-Policy Distillation
% Paragraph 3: OPD Problem 
% Paragraph 4: Method 
% Parargraph 5: Summarization of Results 




% Recent advances in Large Language Models~(LLMs) combined with reinforcement learning~(RL) have produced models with strong performance on complex reasoning tasks~\cite{guo2025deepseek, jaech2024openai, yang2025qwen3}, including mathematical, logical, and multi-step problem solving.
% These advances have sparked growing interest in how such reasoning capabilities can be transferred to smaller, more efficient models with lower inference cost and improved deployability.
% % At the same time, they have motivated growing interest in how such reasoning abilities can be effectively transferred to smaller and more efficient models that offer advantages in inference cost and deployability.

% Knowledge distillation~\cite{hinton2015distilling} is a promising approach for this transfer. Traditionally, distillation relies on off-policy teacher data, where students are trained using supervised loss or forward KL divergence~\citep{kim2016sequence}; However, this approach introduces a distribution mismatch between the sequences observed during training and those generated by the student model at inference time. 
% Recently, on-policy distillation has gained attention, where the student generates samples that are corrected by the teacher via reverse KL divergence~\citep{gu2023minillm, agarwal2024policy, lu2025onpolicydistillation}. 
% Furthermore, incorporating RL into on-policy distillation leverages the strengths of both paradigms, resulting in greater efficiency by requiring fewer rollouts to improve student performance. As illustrated in \autoref{fig:intro_example}, on-policy distillation is highly effective, enabling student models to match the performance of RLVR-trained models within only a few training steps.

% In the context of language models, knowledge distillation trains a student model to match the behavior of a teacher model using objectives defined over token-level predictions. These objectives include supervised fine-tuning on teacher-generated outputs~\cite{kim2016sequence} and distribution matching objectives using logit-based KL~\cite{kullback1951information} divergence.
% Such knowledge distillation is typically performed by minimizing the forward KL divergence between the teacher and student distributions over off-policy data generated by the teacher. However, this approach introduces a distribution mismatch between the sequences observed during training and those generated by the student model at inference time. To address this, on-policy distillation approaches~\cite{gu2023minillm, agarwal2024policy} have been proposed that train the student on self-generated sequences and encourage it to focus on the major modes of the teacher distribution via reverse KL divergence.

% In this paper, we focus on a recently introduced on-policy distillation approach~\cite{lu2025onpolicydistillation} that optimizes student models using dense rewards based on token-level reverse KL divergence over student-generated trajectories. As illustrated in \autoref{fig:intro_example}, on-policy distillation is highly effective, enabling student models to match the performance of RLVR-trained models within only a few training steps.

% However, while on-policy distillation using a reverse KL–based reward effectively captures the dominant modes of the teacher, it tends to reduce student diversity after training and provides unstable learning signals at tokens where the teacher distribution has high entropy. This, in turn, limits the preservation of the teacher’s distribution when probability mass is allocated across multiple tokens---which are usually critical parts, as high-entropy tokens often represent key steps in reasoning~\cite{wang2025beyond}.

% In this paper, we analyze the entropy distributions of generated tokens from the student and teacher after on-policy distillation and the stability of the reward. While on-policy distillation using a reverse KL–based reward effectively captures the dominant modes of the teacher, it tends to reduce student diversity after training and provides unstable learning signals at tokens where the teacher distribution has high entropy. This, in turn, limits the preservation of the teacher’s distribution when probability mass is allocated across multiple tokens.

% % (TW version adapted from YJ's version)
% Motivated by this observation, we propose \methodfullname~(\method), which incorporates the forward KL divergence into the on-policy distillation objective, significantly mitigating these issues. Our experiments clarify the role of entropy in on-policy distillation and reveals that forward KL provides stronger and more stable guidance precisely in high-entropy regions, where reverse KL alone yields weak or unstable signals. As a result, the student is encouraged to preserve probability mass across multiple plausible continuations and more faithfully reflect the teacher’s underlying structure.

% {Yongjin's version}Based on these observations, we propose a new distillation method~(METHOD NAME) that integrates forward KL into the reverse KL framework, preserving the training efficiency of on-policy distillation while capturing the uncertainty and diversity of the teacher distribution. 
% Specifically, while we apply a reverse KL loss across all tokens, we introduce additional guidance using forward KL to better reflect the teacher's structure at high-entropy tokens. 
% This stabilizes learning signals in uncertain regions and prevents the student from converging to a deterministic policy, thereby preserving the teacher's intended diversity. 



% \yj{Summarize the results here} % results/analysis
% Empirically, this guidance maintains student output entropy at a level x\% higher than standard on-policy distillation. Furthermore, it improves [Metric Name]@8 accuracy by x\% and pass@k accuracy by x\% across a range of reasoning benchmarks.

% Based on these observations, in this paper, we propose a new distillation method that utilizes both forward KL and reverse KL, thus preserving the efficient training structure of on-policy distillation while incorporating the uncertainty and diversity of the teacher distribution. 
% Specifically, we apply reverse KL loss for every token. 
% But we introduce additional guidance that reflects the structure of the teacher distribution at high-entropy tokens with forward KL. 
% This stabilizes learning signals in such regions and prevents the student from converging to a deterministic policy, thereby better preserving the diversity intended by the teacher. 
% Empirically, this guidance maintains the output entropy of the student model at a level that is x\% higher than that of standard on-policy distillation, indicating improved preservation of the uncertainty and diversity of the teacher distribution. In addition, it improves average @8 accuracy by x\% and pass@k accuracy by x\% across a range of reasoning benchmarks.

% Our key contributions are: ... 

%These observations raise the following question:
% \begin{quote}
% \emph{Does on-policy distillation adequately utilize the teacher’s distributional information during student learning?}
% \end{quote}

% in progress
% In particular, we consider whether incorporating additional supervision from the teacher 

% To address this limitation, we argue that utilizing forward KL in an adaptive manner is necessary to cover the diverse modes of the teacher distribution. Specifically, we ...